\question[3]
Im Telefonnetz erkennt man an der Vorwahl, ob ein Anruf lokal (Ortsnetz)
oder in ein anderes Netz (z.\,B. anderes Bundesland oder Ausland) geht.
Erkläre die Parallele zwischen der Vorwahl im Telefonnetz und der Subnetzmaske in IP-Netzen. Wozu dient jeweils dieser "Trennmechanismus"?

\begin{solutionbox}{6cm}
Die Vorwahl unterteilt eine Telefonnummer in einen Teil, der
das Ortsnetz (Region) angibt, und einen Teil, der den
Teilnehmer innerhalb dieses Ortsnetzes identifiziert.

Die Subnetzmaske unterteilt eine IP-Adresse ähnlich in einen
Netzteil (welches lokale Netz) und einen Geräteteil (welches
Gerät in diesem Netz).

In beiden Fällen dient der Trennmechanismus dazu, schnell zu
erkennen, ob sich Sender und Empfänger im selben lokalen
Bereich befinden oder ob die Nachricht/Verbindung über eine
Vermittlungsstelle (Router bzw. Vermittlung im Telefonnetz) in
ein anderes Netz weitergeleitet werden muss.
\end{solutionbox}